%% title: Appendix: margin, nuisance augmentation, head scaling, recipe ladder, noise floors
%% target: Sections/Appendix.tex
%% anchor: The complete decoder output sequence is obtained after $N$ layers followed by a final normalization block: $\mathcal{D}(T, e) = \mathrm{RMSNorm}_{\text{final}}\left( T^{(N)} \right)$.
%% action: after
%% why: Results that support the main text but are not part of its argument: why the margin hurts the coarse ranks, how much of the gap is nameable nuisance, why the classifier matrix limits scale, how recipe L was chosen, and the retraining spread per regime in the current convention.
%% source: paper/DRAFT.md §3.1, §4.7, §4.8, §4.12; journal/2026-10-22-rescoring-in-the-mini-metrics-convention.md (repeats); data/rescore/*-grouped.csv (A1 0.3636, B1 0.4086 full trap; repeats base/A5, G3/G3b, R5/R5b, P5/P5b)
%% check: Paste at the end of the appendix. | The recipe ladder documents recipe L only; drop it if recipe L is not kept. | Test-fold numbers are unchanged in the current convention; trap numbers are rescored (grouped calibration). | The ToL-200M floor numbers (884,662 to 203,878 species) assume a 50-image floor, which the merged paper drops for training; keep them as a statement about scale.
%% status: ready
\section{The angular margin and marginalisation}\label{appendix:margin}
Adding an ArcFace margin to the flat head costs 1.00 point at species, 1.15 at genus and 1.11 at family on the test fold, and 0.95, 1.21 and 1.13 on a backbone ten times larger.
The coarse ranks lose more than the species rank they are summed from.
A genus posterior $P(g) = \sum_{s \in g} P(s)$ depends on how probability is spread over the species of a genus, not only on which ranks first; a margin trains against a harder target than the label and distorts that spread.
Supervising the marginals does the opposite (species unchanged, genus $+0.27$, family $+0.39$), so a model that marginalises and uses a margin should also supervise the marginals.

\section{Nameable nuisance}\label{appendix:nuisance}
Adding motion blur, low light and JPEG quantisation to training (B1) raises species macro-$\text{F}_1$ on the 47,905 trap crops from 0.364 to 0.409 (calibrated, grouped) for 0.36 points on the test fold, at no inference cost.
Each transform encodes a guess about what differs between the domains, so the gain measures the part of the gap that can be named in advance; it is worth little once the classifier is adapted to trap images (\cref{sec:results_classifier}).

\section{The classifier matrix at scale}\label{appendix:scaling}
At an embedding size of 1,280 a prototype matrix for the roughly one million species of a global taxonomy takes 15.4\,GB in fp32 with its optimiser state.
For the margin head, class centroids replace the trained matrix at inference for 0.29 points (3.7 for the plain head), and single means beat $k$-means centroids and medoids.
But the matrix is not low-rank (rank 1,035 of 1,280 for 90\,\% of its energy; rank 128 costs 3.1 points), and training without it fails: sampled softmax degrades with no plateau, taxonomy-aware negatives recover 26\,\% of that loss, and a head of running class centroids with no matrix costs 4.6 points.
Centroids can replace a trained matrix at inference but not substitute for training one.
A 50-image floor reduces TreeOfLife-200M from 884,662 to 203,878 species while removing about 7\,\% of images, which brings the footprint to 3.1\,GB.

\section{How recipe L was chosen}\label{appendix:ladder}
Recipe L was built by single-factor ablation (species macro-$\text{F}_1$ on the test fold): Muon with a flat cosine schedule and heavy augmentation, 0.8297; warm-up, lighter augmentation and looser gradient clipping, 0.8769; a one-cycle schedule, 0.8887; square-root class oversampling, 0.9148; 10 instead of 5 epochs without oversampling, 0.8976.
The second step bundles three changes and cannot be attributed.
The optimiser is in every row; schedule, augmentation and sampler moved the number.
Logit adjustment reached 0.9031 and lowered genus and family: one temperature shared by three label distributions (12,041, 4,333 and 102 classes) suits none of them, while oversampling changes which images are drawn and acts the same at every rank.
The 10-epoch run was still improving, so the 5-epoch budget is a comparison convention, not a converged optimum.

\section{Retraining spread}\label{appendix:noise}
\Cref{tab:noise} gives the difference between two identical trainings, by training regime.
It varies more than tenfold on one benchmark, so a noise floor belongs to the training procedure as well as to the evaluation set.

\begin{table}[htb]
    \centering
    \caption{Absolute difference in species macro-$\text{F}_1$ between two identical trainings (recipe L), calibrated on 10\,\% of the (trap, night) groups. On the test fold, end-to-end repeats differ by 0.0000.}
    \label{tab:noise}
    \begin{tabular}{lcc}
        \toprule
        regime & probe & held-out species \\
        \midrule
        end-to-end, 20\,M, 5 epochs & 0.014 & 0.015 \\
        end-to-end self-training, no gate, 20\,M & 0.011 & 0.024 \\
        frozen-backbone stage, 198\,M, 2 epochs & 0.001 & 0.022 \\
        unfrozen adaptation, 303\,M, 2 epochs & 0.004 & 0.026 \\
        \bottomrule
    \end{tabular}
\end{table}
